\subsection{A guard bound using the actual stopping depth}
\label{sec:stopped-guard}
The complex network is still evaluated exactly until a complete normalized
butterfly layer is finished. Write $W=W_{\rm c}$ and $s=s_{\rm c}$ for that
network, and retain
\[
 E=64(W+m+1)^3,\quad B=s+E,\quad C_0=128mB^2.
\]
At $h=50$ the exact values are
\[
 m=125000,\quad W=369474390296480000,\quad
 s=46184298777878576000000,
\]
so $m\ge3$ and $2\le s<m^5$ by integer comparison. We will use $C_1=2$ and
$\Delta=\lceil C_0d^2\rceil$ for every rational $9/10\le\beta<1$.

Measure dependency depth by complex additions, subtractions, halvings, sign
changes, and multiplication by $i$. A chain follows an operand into a later
result. Concatenating every scalar gate and recursive edge call is a valid
upper bound, even though most do not belong to one dependency chain.
In the retained complex motif each wire touches at most twelve scalar gates,
so there are at most $12W$ gates. A gate output sums at most $W$ terms,
with at most a halving and an addition per term. Evaluating all outputs from
saved inputs uses at most $2W^2$ operations per gate. Inverse-child corrections
cost at most $4s$, and source, sink, and bank corrections at most $4W+4$.
Thus their total is bounded by $24W^3+4s+4W+4<E$.
The new bit motif performs exact permutations of coefficient encodings; it
does not introduce coefficient arithmetic into this depth count.

The individual $C$-kernel formula uses at most eight elementary operations.
Consequently the same conservative recurrence remains valid:
\[
 A(e)\le sA(e/m)+E\quad\hbox{at an internal node},\qquad
 A(e)\le8e\quad(e<d^\beta).
\]
For a root $e\le d$, the actual number $j$ of internal levels is at most
$(1-\beta)\log_m d+1$. If $j>0$, this follows from the last internal size
$e/m^{j-1}\ge d^\beta$; for $j=0$ it is immediate. Hence
\[
 s^j\le s\,d^{5(1-\beta)},\qquad
 A(e)\le8d^\beta s^j+E\sum_{i<j}s^i
       \le s(8+E)d^{5-4\beta}
       \le9B^2d^{7/5}.
\]
Here $s\ge2$, $d\ge1$, $\beta\ge9/10$, and
$s(8+E)\le9B^2$; the estimate also covers a root that is already a leaf.
Keeping $E$ additive is essential: it is a fixed cost at a node, not another
branching factor.

The base-$m$ partition has at most $(m-1)(1+\lfloor\log_m d\rfloor)$
pieces. Since $m\ge3$ and $\log d\le\sqrt d$ for $d\ge1$, this is at most
$2m\sqrt d$. Concatenating their depths therefore costs at most
$18mB^2d^{19/10}$. Preprocessing uses at most $d$ individual kernels, costing
$8d$. The two outer diagonals and the scalar $b^D$ contribute at most
$t+9$ further units, where $t=\lfloor D/2\rfloor$; a shift by $t$ is charged
$t$ denominator bits. Together $8d+t+9\le18d$. Thus the total depth and
shift charge are bounded by
\[
 18mB^2d^{19/10}+18d\le36mB^2d^2<C_0d^2.
\]
As in the original induction, an input in $2^{-p}\mathbb Z[i]$ with modulus
at most one has denominator exponent at most $p+\Delta$ and magnitude at most
$2^\Delta$ throughout. Reserving $p+\Delta$ fractional bits and $\Delta+3$
integer and sign bits per real component therefore represents all values
exactly. Address-swap routines restore complete encodings before coefficient
arithmetic resumes and do not change this argument. If $2\epsilon<1$, then
$\Delta=o(p)$ and the full coefficient width remains $O(p)$.
